\documentclass[11pt]{article}

\usepackage{hyperref}
\usepackage{xurl}
\usepackage{float}
\usepackage{array}
\usepackage[margin=1in]{geometry}
\usepackage[tagged, highstructure]{accessibility}
\usepackage{bookmark}
\usepackage{enumitem}

\hypersetup{
    colorlinks=false,
    pdfborder={0 0 0},
    pdftitle={Sample Threat Model: Staying Safe When Leaving a Controlling Partner},
    pdfauthor={Adrianna Holden-Gouveia},
    pdfsubject={Information Security - Sample threat model for the Encrypted Communications and Threat Models Lab},
    pdflang={en-US},
    pdfkeywords={threat model, sample, risk register, tech safety, location tracking, stalkerware}
}

\bookmarksetup{
    numbered,
    open,
}

\setlist{nosep}

\title{Threat Model: Staying Safe When Leaving a Controlling Partner}
\author{SAMPLE ONLY. Fictional person. Not for submission.}
\date{\vspace{-5ex}}

\begin{document}

\maketitle

\section*{Trigger Warning}
\textbf{This sample talks about domestic abuse, stalking, and controlling relationships.} You don't need to read it to do the lab. The other sample threat model, about a streamer, shows all the same parts. If this topic is hard for you, skip this one, and check with me if you'd like another option.

If any of this is happening to you, you're not alone. The National Domestic Violence Hotline is free and confidential: call 1-800-799-7233, text START to 88788, or chat at thehotline.org. Use a device the other person can't get to.

\section*{Read This First}
\begin{itemize}
    \item \textbf{This is a sample.} It shows what a threat model can look like in industry. Sam is a made-up person, and the details are invented.
    \item \textbf{Don't copy it.} Copied wording, risks, or controls get no credit. This scenario isn't one of the lab's choices, so you can't use it.
    \item \textbf{Your lab has its own requirements.} Follow the lab instructions, not this layout. Your lab needs EFF's six questions, at least 10 risks, a plan, and at least 5 sources, each with a CRAAP check, a screenshot, and a quote. This sample has no sources section like that.
    \item \textbf{Not professional advice.} A real safety plan should be made with a trained advocate. Some steps here, like turning off location sharing, can alert the other person, so an advocate can help you decide what to do, and when.
\end{itemize}

\section*{Accessibility Notice}
This document is also available in HTML format at:

Link: \href{https://aholdengouveia.name/InfoSec/labs/threatmodel-leaving.html}{\textbf{\nolinkurl{https://aholdengouveia.name/InfoSec/labs/threatmodel-leaving.html}}}

\section{Document Control}
\begin{table}[H]
\caption{Document control}
\begin{tabular}{|m{0.3\textwidth}|m{0.6\textwidth}|}
\hline
\textbf{Field} & \textbf{Value} \\
\hline
Document ID & TM-SAMPLE-02 \\
\hline
Version & 1.0 \\
\hline
Date & October 15, 2026 \\
\hline
System owner & Sam (fictional) \\
\hline
Prepared by & Adrianna Holden-Gouveia, for CIS 115 \\
\hline
Classification & SAMPLE. Not for submission. \\
\hline
Next review & Weekly while Sam is moving, then every 3 months
\end{tabular}
\end{table}

\section{Summary}
Sam is leaving a controlling partner they lived with for 3 years, and is moving in with a friend. They share a phone plan, a laptop, and a car with a connected app, and the partner knows several of Sam's passwords. The biggest risks are the partner finding Sam's new location through their phone, a hidden tracker, or the car, and reading Sam's messages through shared accounts. The most important protections are a new phone in Sam's own name, securing email first, turning off location sharing, and checking for trackers, all planned with an advocate. After these protections, the remaining risk is Medium, mostly because the partner may still try to reach Sam through friends and family.

\section{Scope}
\textbf{In scope:} Sam's phone, accounts, messages and calls, location, and the shared laptop, phone plan, and car.

\textbf{Out of scope:} legal protection, such as a protective order, housing, money, and physical safety plans. An advocate and a lawyer cover these.

\section{Assets}
\begin{table}[H]
\caption{Assets to protect}
\begin{tabular}{|m{0.3\textwidth}|m{0.6\textwidth}|}
\hline
\textbf{Asset} & \textbf{Why it matters} \\
\hline
Sam's new address and location & If the partner finds it, Sam could be harassed or hurt \\
\hline
Messages and calls with family, the lawyer, and the advocate & They show Sam's plans \\
\hline
Email account & It can reset the password on almost every other account \\
\hline
Bank and other accounts & The partner could take money or lock Sam out \\
\hline
Phone and phone number & The phone holds every account, and the number receives some login codes
\end{tabular}
\end{table}

\section{Threat Actors}
\begin{table}[H]
\caption{Who might attack, and why}
\begin{tabular}{|m{0.25\textwidth}|m{0.3\textwidth}|m{0.35\textwidth}|}
\hline
\textbf{Threat actor} & \textbf{What they want} & \textbf{What they can do} \\
\hline
The ex-partner & To find Sam, and keep control & Knows Sam's passwords and habits, has had access to Sam's phone, owns the phone plan \\
\hline
People the partner asks for help & To ``help'' the partner reach Sam & Ask Sam's friends and family where Sam is
\end{tabular}
\end{table}

\section{Assumptions}
\begin{itemize}
    \item The partner has average tech skills, and is willing to use apps and trackers that are easy to buy.
    \item Sam's new phone hasn't been touched by the partner.
    \item The friend Sam is staying with can be trusted.
\end{itemize}
If any of these stop being true, update this threat model, and talk to the advocate.

\section{Entry Points}
These are the ways the partner could reach Sam's assets:
\begin{itemize}
    \item Location sharing settings and family locator apps on Sam's phone
    \item Any app the partner may have installed on Sam's phone
    \item Sam's bag, car, and belongings, where a tracker could be hidden
    \item The car's connected app
    \item Accounts where the partner knows the password, or is still logged in
    \item The shared phone plan, laptop, and cloud accounts
    \item Sam's friends, family, and social media
\end{itemize}

\section{Risk Rating}
Each threat gets a likelihood (how likely it is to happen) and an impact (how bad it would be), from Low to High. The overall risk comes from this table.
\begin{table}[H]
\caption{How likelihood and impact combine into an overall risk}
\begin{tabular}{|m{0.22\textwidth}|m{0.2\textwidth}|m{0.2\textwidth}|m{0.2\textwidth}|}
\hline
\textbf{Likelihood} & \textbf{Low impact} & \textbf{Medium impact} & \textbf{High impact} \\
\hline
High & Medium & High & High \\
\hline
Medium & Low & Medium & High \\
\hline
Low & Low & Low & Medium
\end{tabular}
\end{table}

\section{Threat Register}
Ordered from highest to lowest risk.
\begin{table}[H]
\caption{Threat register}
\begin{tabular}{|m{0.06\textwidth}|m{0.46\textwidth}|m{0.11\textwidth}|m{0.11\textwidth}|m{0.1\textwidth}|}
\hline
\textbf{ID} & \textbf{Threat} & \textbf{Likeli\-hood} & \textbf{Impact} & \textbf{Risk} \\
\hline
T1 & Location sharing on Sam's phone, such as Find My, Google Maps, or a family locator app, shows where Sam is & High & High & High \\
\hline
T2 & Stalkerware, an app that secretly reports texts, calls, and location, is on Sam's phone & Medium & High & High \\
\hline
T3 & A hidden tracker, such as an AirTag, is in Sam's bag or car & Medium & High & High \\
\hline
T4 & The car's connected app shows where the car is parked & Medium & High & High \\
\hline
T5 & The partner logs in to Sam's email with a password they know, and resets Sam's other passwords & High & High & High \\
\hline
T6 & As owner of the phone plan, the partner sees call and text logs, or moves Sam's number to another SIM card & High & Medium & High \\
\hline
T7 & Shared cloud accounts sync Sam's photos and messages to the shared laptop or a tablet, and photos can include where they were taken & Medium & Medium & Medium \\
\hline
T8 & Friends post or tag photos of Sam, or apps with location permission show where Sam goes & Medium & Medium & Medium \\
\hline
T9 & The shared laptop's saved passwords and browser history show Sam's searches for apartments, lawyers, or advocates & Medium & Medium & Medium \\
\hline
T10 & The partner contacts Sam's friends and family, pretending it's an emergency, to find out where Sam is & Medium & Medium & Medium
\end{tabular}
\end{table}

\section{Protections}
Plan the timing of these with an advocate. Changing settings on the old phone or accounts can alert the partner, so some steps may need to wait until Sam is somewhere safe.
\begin{table}[H]
\caption{Protections, and the threats each one reduces}
\begin{tabular}{|m{0.06\textwidth}|m{0.6\textwidth}|m{0.12\textwidth}|m{0.1\textwidth}|}
\hline
\textbf{ID} & \textbf{Protection} & \textbf{Threats} & \textbf{When} \\
\hline
C1 & Get a new phone on a new plan, in Sam's own name, with a new number and a strong passcode. Don't sign in to old shared accounts on it. & T1, T2, T6 & Before leaving \\
\hline
C2 & From the new phone, change the email password first, turn on two-factor authentication with an authenticator app, and sign out all other devices & T5 & Before leaving \\
\hline
C3 & Change passwords for banking, social media, and cloud accounts, and stop using shared accounts & T5, T7 & After leaving \\
\hline
C4 & Ask an advocate to help check the old phone for stalkerware, and decide when to stop using it & T1, T2 & With an advocate \\
\hline
C5 & On the new phone, keep location sharing off, and turn off location, camera, and microphone for apps that don't need them & T1, T8 & Ongoing \\
\hline
C6 & Keep unknown tracker alerts turned on. Both iPhone and Android can warn about an unknown AirTag, and Android can also run a manual scan. Check bags and the car. If a tracker is found, talk to an advocate or the police before removing it. & T3 & Ongoing \\
\hline
C7 & Remove Sam from the car's app, or ask the lawyer about it if the car is in the partner's name & T4 & With the lawyer \\
\hline
C8 & Stop using the shared laptop. Use a library computer or the new phone instead. & T9 & After leaving \\
\hline
C9 & Ask friends and family not to post, tag, or share Sam's location or new number, even if the partner says it's an emergency & T8, T10 & Before leaving
\end{tabular}
\end{table}

\subsection*{Communications}
\textbf{Messages and calls with family, the lawyer, and the advocate:} Signal on the new phone, which is end-to-end encrypted, so only the people in the chat can read the messages. Turn on disappearing messages, and Registration Lock, which uses a Signal PIN, so no one else can sign up for Signal with Sam's number. \textbf{Regular calls and texts:} these aren't end-to-end encrypted, and the phone company keeps records of them, so use them only for things that aren't sensitive.

\subsection*{VPN decision}
A VPN isn't a priority here. The partner isn't watching Sam's network. They have access to Sam's accounts and devices. A VPN hides which websites you visit from the network you're on, but it doesn't stop location sharing, stalkerware, trackers, or someone who knows your password. If Sam ever has to use the old home Wi-Fi, a well-known, paid VPN would hide Sam's browsing from that router.

\section{Accepted Risk}
After the protections, the remaining risk is Medium:
\begin{itemize}
    \item The partner may still find out things through mutual friends (T10). Sam can lower this risk, but can't control what other people say.
    \item The partner may guess Sam's location from Sam's job or usual places. The advocate and lawyer can help with this, and it's outside this model.
\end{itemize}

\section{Not Covered}
This model only covers technology. Physical safety, legal protection, and housing need a safety plan made with an advocate.

\section{When to Review}
Review this threat model every week while Sam is moving, then every 3 months, and right away if:
\begin{itemize}
    \item The partner seems to know something they shouldn't, such as where Sam was
    \item A tracker alert goes off, or an account shows a login Sam doesn't recognize
    \item Sam moves, changes jobs, or gets a new phone
\end{itemize}

\section{References}
The National Domestic Violence Hotline
Link: \href{https://www.thehotline.org/}{\textbf{\nolinkurl{https://www.thehotline.org/}}}

National Network to End Domestic Violence: Technology Safety
Link: \href{https://www.techsafety.org/}{\textbf{\nolinkurl{https://www.techsafety.org/}}}

EFF Surveillance Self-Defense: Your Security Plan
Link: \href{https://ssd.eff.org/module/your-security-plan}{\textbf{\nolinkurl{https://ssd.eff.org/module/your-security-plan}}}

The lab this sample goes with
Link: \href{https://www.aholdengouveia.name/InfoSec/labs/enccomm.html}{\textbf{\nolinkurl{https://www.aholdengouveia.name/InfoSec/labs/enccomm.html}}}

\end{document}
